%! Populations, Samples, Parameters, and Statistics — Unit 1, Lesson 1.2 — Guided Notes & Practice (Answer Key)
% PILOT SHAPE (2026-09-06): modeled on AP Statistics lesson 1.4 minus its AP Practice page.
% THE in-class document. The notes carry the whole gradual release, closed by a SPOKEN debrief.
% They run SHORT on purpose: the period ends with students starting the homework, which is
% where the individual practice lives.
%   I DO   : the objective box, the vocabulary box, and notes sections 1-2.          (~20 min)
%   WE DO  : the practicebox ("Guided Practice") — the Riverbend job survey, together. (~15 min)
%   YOU DO : nothing on this page — ../homework/ IS the individual practice, started in the
%            last ten minutes of the period.
% No hookbox (the hook is a slide and lives in the plan), no closing practice set, no extension.
% Vocabulary is GIVEN here, not discovered. THE CRUX is the second half of section 2: a
% statistic that differs from another statistic is not a mistake — the parameter is one fixed
% number and statistics vary by sample (PS.DC.1d-e). The second half of section 1 carries the
% second trap: a percent computed from records of EVERYONE is a parameter, not a statistic.
%
% THE DATA
%   Lakeside Farmers Market: 800 Saturday shoppers; 50 asked, 26 drove -> 52%; 0.52 x 800 = 416
%   Three volunteers, a different 50 shoppers each: Ana 26 (52%), Ben 24 (48%), Cleo 29 (58%)
%   District records: 61% of all 900 Riverbend students ride a bus  (a PARAMETER — the trap row)
%   Guided Practice, Riverbend HS: 900 students; 50 asked, 20 have a job -> 40%; 0.40 x 900 = 360;
%     a second 50 asked, 23 have one -> 46%
%
% PAGE LOCKSTEP with notes_key/: the key is GENERATED from this file — every \blank{} becomes
% an \ans{}, every \par\writelines{n} becomes n \ansline{} calls, every \termblank{} a
% \vocabans{}{}. Every \begin{work} block is byte-identical. Blank widths are sized to their
% answers so the two files wrap alike.
\documentclass[12pt]{article}
\usepackage{saar-article}
\usepackage{saar-key}   % pulls in -boxes; do NOT also load -boxes
\newcommand{\TallMath}[1]{$\displaystyle #1\rule[-1.4em]{0pt}{3.2em}$}

\begin{document}

\pageheader{Unit 1, Lesson 1.2}{Guided Notes \& Practice --- Answer Key}
% NAMESTRIP: no \namedateperiod here — mirrors the blank. NO teachernote here —
% teacher prose lives in the lesson plan.

\vspace{0.15in}

\begin{objectivebox}
\small
By the end of these notes I will be able to\dots
\begin{enumerate}[leftmargin=1.5em, itemsep=2pt, topsep=2pt]
  \item name the \textbf{population} and the \textbf{sample} in a study, and give the sample
        size;
  \item tell a \textbf{parameter} from a \textbf{statistic} by asking \emph{who the number
        describes} --- never by what the number looks like;
  \item name the \textbf{constraint} --- time, cost, access, destruction, or change --- that
        keeps a study from being a \textbf{census};
  \item explain why two honest samples give two different statistics for one unchanging
        parameter (\textbf{sampling variability}), and say what a sample does --- and does
        \emph{not} --- let a report claim.
\end{enumerate}
\end{objectivebox}

\boxguard
\begin{vocabbox}
\small
Fill in each term \textbf{as we name it} in the notes below.
\vocabans{Population}{The entire group the question is about --- every individual you would like the conclusion to describe. Decided by the question, not by who is available.}
\vocabans{Sample}{The part of the population you actually collected data from. The number of individuals in it is the sample size, $n$.}
\vocabans{Parameter}{A number that describes the population. Usually unknown --- and there is exactly one of it, whether or not anyone has ever computed it.}
\vocabans{Statistic}{A number computed from the sample. Known, because you collected the data --- and it changes from one sample to the next.}
\vocabans{Sampling variability}{The change in a statistic from one sample to the next, even though the population did not change. Not a mistake --- it is what samples do.}
\end{vocabbox}

% ── I DO: section 1 ──────────────────────────────────────────────────────────
% Guarded at 12 in BOTH files (the vocab rows match in height, so no float-early risk).
\boxguard[12]
\begin{notesbox}{1. Two Groups --- and the Two Kinds of Numbers They Produce}
\small
Last night's homework asked $50$ shoppers and then talked about all $800$. Those two groups
have names, and so do the numbers computed from them.

\vspace{3pt}
The \textbf{population} is the entire group the question is about --- every individual you
would like the conclusion to describe. At Lakeside, that is all \ans{800} Saturday
shoppers. The \textbf{sample} is the part of the population you actually collected data from
--- the \ans{50} shoppers the manager asked. The number of individuals in the sample is
the \textbf{sample size}, $n = $ \ans{50}.

\vspace{3pt}
Name the two groups in each study. \emph{The sample is always chosen out of the population
--- so before you can name a sample, you have to know what the population is.}
\nopagebreak

\begin{center}
\renewcommand{\arraystretch}{1.4}
\begin{tabularx}{\linewidth}{@{} X p{3.3cm} p{2.8cm} c @{}}
\toprule
\textbf{Study} & \textbf{Population} & \textbf{Sample} & \textbf{Size} \\
\midrule
The manager asks $50$ of the market's $800$ Saturday shoppers how they traveled.
  & \ans{the 800 Saturday shoppers} & \ans{the 50 asked} & \ans{50} \\
The principal asks $45$ of Riverbend's $900$ students how many hours a week they work.
  & \ans{Riverbend's 900 students} & \ans{the 45 asked} & \ans{45} \\
A cereal plant weighs $30$ of the $12{,}000$ boxes it filled on Tuesday.
  & \ans{the 12,000 boxes filled Tuesday} & \ans{the 30 weighed} & \ans{30} \\
\bottomrule
\end{tabularx}
\end{center}

\vspace{3pt}
Collecting data from \emph{every} individual in the population is called a \ans{census}. The
U.S. Census is one; almost nothing else you meet in this course will be.

\vspace{10pt}
\textbf{\textcolor{cerulean}{Parameter and Statistic.}}
Both groups produce numbers, and the two kinds are not interchangeable.

\vspace{3pt}
A number that describes the \textbf{population} is a \ans{parameter}. \quad
A number that describes the \textbf{sample} is a \ans{statistic}.

\vspace{2pt}
\textit{Memory hook:} \textbf{p}opulation goes with \textbf{p}arameter; \textbf{s}ample goes
with \textbf{s}tatistic. To sort a number, ask one question: \textbf{who does it describe?}
\nopagebreak

\begin{center}
\renewcommand{\arraystretch}{1.45}
\begin{tabularx}{\linewidth}{@{} X p{2.6cm} p{2.6cm} @{}}
\toprule
\textbf{The number} & \textbf{Describes\dots} & \textbf{So it is a\dots} \\
\midrule
$52\%$ of the $50$ shoppers asked drove to the market.
  & \ans{the sample} & \ans{statistic} \\
The market had $800$ shoppers on Saturday morning.
  & \ans{the population} & \ans{parameter} \\
The $8$ shoppers asked traveled a mean of $4$ miles.
  & \ans{the sample} & \ans{statistic} \\
District records show $61\%$ of Riverbend's $900$ students ride a bus.
  & \ans{the population} & \ans{parameter} \\
\bottomrule
\end{tabularx}
\end{center}

\vspace{4pt}
\begin{center}
\fcolorbox{redacc}{redbg}{\parbox{0.94\linewidth}{\small
\textbf{Careful --- a percent is not automatically a statistic.} The $61\%$ was computed from
the records of \ans{all} $900$ students --- every individual in the population. It
describes the \ans{population}, so it is a \ans{parameter}, even though it is a percent and even
though somebody had to compute it. The question is never \emph{what does the number look
like}; it is \emph{\ans{who} does it describe}.}}
\end{center}

\vspace{4pt}
Which kind do you almost always \emph{know}? The \ans{statistic}. \quad Which kind do you
actually \emph{want}? The \ans{parameter}. That gap is the whole reason this course exists: you
use the number you have to \textbf{estimate} the number you want.

\vspace{2pt}
{\small\textit{Notation you will meet later in the course: a population mean is written $\mu$
and a sample mean $\bar{x}$; a population percent is written $p$ and a sample percent
$\hat{p}$. You do not need these symbols today.}}
\end{notesbox}

% ── I DO: section 2 — constraints, then the crux ─────────────────────────────
\boxguard[26]
\begin{notesbox}{2. Why We Sample at All --- and What Happens When We Do}
\small
If the manager wants a parameter, why not just ask all $800$? A \textbf{constraint} is
something in the real world that limits how a study can be run. Constraints are the reason a
census is usually off the table --- and they are part of describing a study honestly.
\nopagebreak

\begin{center}
\renewcommand{\arraystretch}{1.25}
\begin{tabularx}{\linewidth}{@{} p{2.6cm} X @{}}
\toprule
\textbf{Constraint} & \textbf{What it means} \\
\midrule
Time        & Reaching everyone would take longer than you have. \\
Cost        & Reaching everyone would cost more than the budget. \\
Access      & Some individuals cannot be found, reached, or made to answer. \\
Destruction & Measuring the individual ruins it, so you cannot measure them all. \\
Change      & The population keeps changing while you collect. \\
\bottomrule
\end{tabularx}
\end{center}

\vspace{3pt}
Name the constraint that stands in the way in each situation.
\nopagebreak

\begin{center}
\renewcommand{\arraystretch}{1.45}
\begin{tabularx}{\linewidth}{@{} X p{2.8cm} @{}}
\toprule
\textbf{Situation} & \textbf{Constraint} \\
\midrule
To check a cereal box's weight you must open it, and $12{,}000$ were filled Tuesday.
  & \ans{destruction} \\
A different set of shoppers comes to the market every Saturday.
  & \ans{change} \\
Surveying all $900$ Riverbend students would take three class periods nobody can spare.
  & \ans{time} \\
Reaching every household in the county would mean mailing $40{,}000$ letters.
  & \ans{cost} \\
Many shoppers pay and leave before anyone can get near them with a clipboard.
  & \ans{access} \\
\bottomrule
\end{tabularx}
\end{center}

\vspace{10pt}
\textbf{\textcolor{cerulean}{The Statistic Moves. The Parameter Does Not.}}
Three volunteers each asked a \emph{different} $50$ shoppers on the \emph{same} Saturday
morning. Here is how many in each group said they drove.
\nopagebreak

\begin{center}
\renewcommand{\arraystretch}{1.4}
\begin{tabularx}{\linewidth}{@{} l c X @{}}
\toprule
\textbf{Volunteer} & \textbf{Drove} & \textbf{Percent of that sample} \\
\midrule
Ana   & 26 & $26/50 = 0.52 = 52\%$ \\
Ben   & 24 & \ans{$24/50 = 0.48 = 48\%$} \\
Cleo  & 29 & \ans{$29/50 = 0.58 = 58\%$} \\
\bottomrule
\end{tabularx}
\end{center}

\vspace{3pt}
Show the work for \textbf{Cleo's} row:

\begin{work}
  \dfrac{29}{50} &= 0.58 \\
                 &= 58\%
\end{work}

\vspace{2pt}
How many different statistics did the three volunteers produce? \ans{3} \qquad
How many true percents are there for all $800$ Saturday shoppers? \ans{1}

\vspace{3pt}
Did one of the volunteers count wrong? \ans{No} \quad Which of the three found the true
percent for all $800$? \ans{none --- each is an estimate}

\vspace{5pt}
\begin{center}
\fcolorbox{cerulean}{frost}{\parbox{0.94\linewidth}{\small
\textbf{A different answer is not a wrong answer.} This change from one sample to the next
is called \ans{sampling variability}. It is not a mistake --- nobody miscounted. It is what samples
\ans{do}. The statistic \ans{moves} from sample to sample; the parameter
\ans{stays put}, because the $800$ shoppers did not change.}}
\end{center}

\vspace{4pt}
Each of the three statistics is an \textbf{estimate} of the one parameter. A larger sample
usually lands \ans{closer} to it, but no sample lands exactly on it every time.

\vspace{6pt}
\textbf{One limit on what you may claim.} Fill in the manager's report sentence from her own
sample, then judge two sentences she is tempted to write instead.

\vspace{2pt}
\textit{``Of the \ans{50} shoppers we asked, \ans{52}\% drove. If our sample
represents the market, that is about \ans{416} of the \ans{800} shoppers who come on
a Saturday.''}
\nopagebreak

\begin{center}
\renewcommand{\arraystretch}{1.5}
\begin{tabularx}{\linewidth}{@{} X p{2cm} >{\raggedright\arraybackslash}p{5.2cm} @{}}
\toprule
\textbf{Claim} & \textbf{Allowed?} & \textbf{Why} \\
\midrule
``Exactly $416$ Saturday shoppers drove to the market.''
  & \ans{No} & \ans{estimate, not a count} \\
``$52\%$ of Lakeside residents drive to the market.''
  & \ans{No} & \ans{wrong population} \\
\bottomrule
\end{tabularx}
\end{center}
\end{notesbox}

% ── WE DO: guided practice ───────────────────────────────────────────────────
\boxguard[24]
\begin{practicebox}
\small
\textbf{The Riverbend Job Survey --- we work this one together.} The Riverbend principal wants
to know what fraction of the school's $900$ students hold a part-time job. She asks $50$
students; $20$ of them have one.

\vspace{3pt}
\textbf{(a)} Name the groups.

Population: \ans{Riverbend's 900 students} \hfill Sample size: \ans{50}

Sample: \ans{the 50 students she asked}

\vspace{3pt}
\textbf{(b)} What percent of the students she asked have a part-time job, and about how many of
Riverbend's $900$ students would that be?

\begin{work}
  \dfrac{20}{50}  &= 0.40 = 40\% \\
  0.40 \times 900 &= 360 \text{ students}
\end{work}

The $40\%$ is a \ans{statistic}. The $360$ is an \ans{estimate} of a \ans{parameter}.

\vspace{3pt}
\textbf{(c)} The same week, the assistant principal asks a \emph{different} $50$ students the
same question. In that group, $23$ have a part-time job.

\begin{work}
  \dfrac{23}{50} &= 0.46 = 46\%
\end{work}

One survey says $40\%$; the other says $46\%$. Did one of them make a mistake? Which survey
found the true percent for all $900$ students? Explain.
\par\ansline{No. Two different samples of 50 give two different statistics.}
\ansline{Neither found it --- there is one parameter, and both estimate it.}

\vspace{3pt}
\textbf{(d)} Name one constraint that kept the principal from asking all $900$ students, and say
which of the five it is.
\par\ansline{Time --- three class periods nobody can spare (also cost or access).}
\end{practicebox}

\end{document}
